\documentclass[11pt, english, letterpaper]{article}
\usepackage[T1]{fontenc}
\usepackage[utf8]{inputenc}
\usepackage{textcomp}
\usepackage{lmodern}
\usepackage{graphicx}
\usepackage{pdfpages}
\usepackage[usenames,dvipsnames,table]{xcolor}
\usepackage[margin=0.9in]{geometry}
\usepackage{babel}
\usepackage{caption}
\usepackage{float}
\usepackage{longtable}
\usepackage{booktabs}
\usepackage{tabularx}
\usepackage[shortlabels]{enumitem}
\usepackage{url}
\usepackage[colorlinks=true, linkcolor=black, urlcolor=black, citecolor=black,
linkbordercolor={white}, anchorcolor=black]{hyperref}
%
% Glyphs the Python exhibits emit that pdflatex does not map by default.
% \ensuremath is used instead of dollar-sign math: Stata would try to expand
% anything after a dollar sign as a global macro.
\DeclareUnicodeCharacter{2265}{\ensuremath{\geq}}
\DeclareUnicodeCharacter{2264}{\ensuremath{\leq}}
\DeclareUnicodeCharacter{00D7}{\ensuremath{\times}}
\DeclareUnicodeCharacter{2212}{--}
\DeclareUnicodeCharacter{2013}{--}
\DeclareUnicodeCharacter{2014}{---}
\DeclareUnicodeCharacter{00B0}{\textdegree}
\DeclareUnicodeCharacter{00B2}{\textsuperscript{2}}
\DeclareUnicodeCharacter{2018}{\textquoteleft}
\DeclareUnicodeCharacter{2019}{\textquoteright}
\DeclareUnicodeCharacter{201C}{\textquotedblleft}
\DeclareUnicodeCharacter{201D}{\textquotedblright}
%
\pdfminorversion=6
\linespread{1.15}
\def\arraystretch{1.25}
\captionsetup{font=small, labelfont=bf, justification=raggedright, singlelinecheck=false}
\setlength{\parskip}{6pt}
\setlength{\parindent}{0pt}
%
% Exhibits live beside this .tex, so the folder compiles anywhere.
\newcommand*{\figTwo}{exhibits}%
%
\begin{document}
\title{Winter Weather Panel: Descriptive Report\\[4pt]
\large PRISM county--month data, 1981--2025}
\author{Wendy Wang}
\date{2026-08-28}
\maketitle
\thispagestyle{empty}
\vspace{-1em}
\begin{center}\begin{minipage}{0.86\textwidth}\small
\textbf{Provenance.} All exhibits generated by
\path{codePYTHON/09_descriptive_weather_full.ipynb}, executed on Kodama against the
complete PRISM county-month panel. Figures are vector PDFs reproduced without
rescaling of content. Tier~1 QA exhibits are indexed in Appendix~B rather than
reproduced.
\par\vspace{6pt}
\textbf{Status.} The warm-anomaly exhibits in Section~\ref{sec:anom} are
\textbf{provisional} pending confirmation of the anomaly definition.
Section~\ref{sec:decisions} lists the choices requiring a decision.
\end{minipage}\end{center}
\newpage
\tableofcontents
\newpage
\section{Coverage and data integrity}
The panel is complete: no duplicate keys, no unmatched rows across the monthly
and derived files, no incomplete county-months, and no QA-field mismatches. The
excluded county-winters are exactly the two boundary winters that cannot be
complete, not unexplained missingness.
\par
\begin{longtable}{@{}lr@{}}
\caption{Coverage and integrity diagnostics}\\
\toprule Metric & Value \\ \midrule \endfirsthead
\toprule Metric & Value \\ \midrule \endhead
\bottomrule \endfoot
dataset & PRISM \\
monthly rows & 1678320 \\
derived rows & 1678320 \\
merged rows & 1678320 \\
counties & 3108 \\
first year & 1981 \\
last year & 2025 \\
duplicate monthly keys & 0 \\
duplicate derived keys & 0 \\
monthly only rows & 0 \\
derived only rows & 0 \\
incomplete months & 0 \\
qa mismatch n days & 0 \\
qa mismatch expected days & 0 \\
qa mismatch is incomplete & 0 \\
qa mismatch dataset types & 0 \\
candidate county winters & 142968 \\
complete county winters & 136752 \\
excluded incomplete county winters & 6216 \\
first complete winter & 1982 \\
last complete winter & 2025 \\
expected calendar years & 1981–2025 \\
observed calendar years & 1981–2025 \\
observed complete winters & 1982–2025 \\
\end{longtable}
\begin{longtable}{@{}lrl@{}}
\caption{Tier 1 QA findings}\\
\toprule Check & Result & Status \\ \midrule \endfirsthead
\toprule Check & Result & Status \\ \midrule \endhead
\bottomrule \endfoot
Duplicate county-year-month keys & 0 & PASS \\
Rows appearing in only one merge source & 0 & PASS \\
Incomplete county-months & 0 & PASS \\
Excluded incomplete county-winters & 6216 & INFO \\
County-variable mean flags (|z| ≥ 5) & 20 & REVIEW \\
\end{longtable}
\clearpage
\section{Long-run winter warming}
\begin{figure}[H]
\centering
\includegraphics[width=\textwidth]{\figTwo/prism_national_long_run_warming.pdf}
\caption{National winter temperature, 1982--2025. Each point is the unweighted
mean across counties, so the estimand is the average county rather than an area-
or population-weighted national average. The fitted slope is shown with its
standard error.}
\end{figure}
State-level trends are reproduced in Appendix~\ref{app:states}.
\clearpage
\section{Where winters have warmed}
\begin{figure}[H]
\centering
\includegraphics[width=\textwidth]{\figTwo/prism_winter_temperature_early_vs_recent.pdf}
\caption{Mean winter temperature in the early (1982--1991) and recent
(2016--2025) periods, on identical colour limits. Both panels are ten-winter
period means, not single-year endpoints.}
\end{figure}
\clearpage
\begin{figure}[H]
\centering
\includegraphics[width=\textwidth]{\figTwo/prism_winter_temperature_change.pdf}
\caption{Change in mean winter temperature between the two periods. The
diverging scale is centred at zero, so the map shows both the breadth of warming
and its regional heterogeneity.}
\end{figure}
\clearpage
\section{Cold extremes}
\begin{figure}[H]
\centering
\includegraphics[width=\textwidth]{\figTwo/prism_extreme_cold_day_count.pdf}
\caption{Mean count of winter days with a daily minimum below
0\textdegree F. This is the extreme-cold measure only; it is not a Winter
Severity Index.}
\end{figure}
\clearpage
\section{Anomalously warm winters}
\label{sec:anom}
\begin{figure}[H]
\centering
\includegraphics[width=\textwidth]{\figTwo/prism_anomalous_warm_winter_frequency_provisional.pdf}
\caption{\textbf{Provisional.} Share of a county's eligible winters classified
as anomalously warm, defined as a detrended residual above 1.5 standard
deviations of that county's own full-sample distribution.}
\end{figure}
\clearpage
\begin{figure}[H]
\centering
\includegraphics[width=\textwidth]{\figTwo/prism_anomalous_warm_winter_time_series_provisional.pdf}
\caption{\textbf{Provisional.} Share of counties classified as anomalously warm
in each winter. Because each county's linear trend is removed, this isolates
episodic shocks from gradual warming.}
\end{figure}
\clearpage
\subsection*{Known properties of the provisional definition}
\begin{enumerate}[leftmargin=*]
\item Look-ahead: each county's reference distribution uses the full 1982–2025 sample, so later winters inform how earlier winters are classified.
\item Detrending removes each county's linear trend, so gradual warming is excluded from the shock measure by construction.
\item A county's mapped frequency partly reflects the skewness and tail shape of its residual distribution, not exposure to warm shocks alone.
\item The 1.5 SD cut-off is a convention, not an estimated threshold.
\item For regression use, compare against a fixed-baseline alternative (e.g. above the county's 1982–2000 mean by 1.5 baseline SD) or use the continuous standardised anomaly.
\end{enumerate}
\clearpage
\section{Distributional shift}
\begin{figure}[H]
\centering
\includegraphics[width=\textwidth]{\figTwo/prism_national_winter_temperature_distributions_by_decade.pdf}
\caption{County-winter mean temperatures pooled by period. Periods contain
unequal numbers of winters, and the pooled spread is dominated by persistent
cross-county geography, which compresses the apparent temporal shift.}
\end{figure}
\clearpage
\begin{figure}[H]
\centering
\includegraphics[width=\textwidth]{\figTwo/prism_national_winter_temperature_deviations_by_decade.pdf}
\caption{The same periods after subtracting each county's own full-sample mean.
Removing persistent geography isolates the within-county shift, which is the
quantity of interest.}
\end{figure}
\clearpage
\section{Robustness: PRISM versus ERA5}
\begin{figure}[H]
\centering
\includegraphics[width=\textwidth]{\figTwo/prism_era5_national_winter_temperature_comparison.pdf}
\caption{National winter temperature from both products. PRISM is the analysis
dataset; ERA5 is a targeted robustness check. The mean offset and correlation are
reported on the figure.}
\end{figure}
\clearpage
\section{Decisions requested}
\label{sec:decisions}
\subsection*{Anomaly definition (regression use)}
\textbf{Current choice.} Detrended county residual > 1.5 SD of the full 1982–2025 sample; minimum 30 winters\par
\textbf{Issue.} Uses future winters to classify earlier ones and removes secular warming by construction; threshold is a convention. Alternatives: fixed 1982–2000 baseline, or the continuous standardised anomaly.\par
\textbf{Status.} CONFIRM — descriptive exhibits are labelled PROVISIONAL
\subsection*{Geographic weighting}
\textbf{Current choice.} Unweighted mean across counties\par
\textbf{Issue.} Estimand is the average county, not an area-, population- or road-weighted national average. A VMT- or road-mile-weighted series may be the relevant exposure weight for the collisions stage.\par
\textbf{Status.} CONFIRM
\subsection*{Comparison periods}
\textbf{Current choice.} 1982–1991 versus 2016–2025\par
\textbf{Issue.} Symmetric ten-winter period means; not 1981-versus-2025 endpoints.\par
\textbf{Status.} CONFIRM
\subsection*{Extreme-cold exhibit}
\textbf{Current choice.} Mean days with TMIN < 0°F; not labelled WSI\par
\textbf{Issue.} Not a full Winter Severity Index.\par
\textbf{Status.} DOCUMENTED
\subsection*{ERA5 robustness exhibit}
\textbf{Current choice.} Targeted national PRISM-ERA5 comparison\par
\textbf{Issue.} PRISM remains the analysis dataset; ERA5 runs warm by a small constant offset.\par
\textbf{Status.} DOCUMENTED
\clearpage
\appendix
\section{State-level winter-temperature trends}
\label{app:states}
Per-state annual means with state-specific linear trends and the standard error
of each slope. Reproduced in full; page count follows the run.
\includepdf[pages=-]{\figTwo/prism_state_long_run_warming.pdf}
\section{Tier 1 QA artefacts (not reproduced)}
The following internal QA exhibits and tables were produced by the same run and
are available on request.
\begin{itemize}[leftmargin=*,itemsep=1pt]
\item \path{figures/weather/tier1/prism_interannual_variability_choropleth.pdf}
\item \path{figures/weather/tier1/prism_national_trend_days_below_freezing_32f.pdf}
\item \path{figures/weather/tier1/prism_national_trend_days_extremely_cold.pdf}
\item \path{figures/weather/tier1/prism_national_trend_freeze_thaw_days.pdf}
\item \path{figures/weather/tier1/prism_county_spaghetti_mean_temp_c.pdf}
\item \path{figures/weather/tier1/prism_state_trends_freeze_thaw_days.pdf}
\item \path{figures/weather/tier1/prism_state_decade_distributions_mean_temp_c.pdf}
\item \path{figures/weather/tier1/prism_state_trends_mean_temp_c.pdf}
\item \path{figures/weather/tier1/prism_state_trends_days_below_freezing_32f.pdf}
\item \path{figures/weather/tier1/prism_national_trend_mean_temp_c.pdf}
\item \path{figures/weather/tier1/prism_state_trends_days_extremely_cold.pdf}
\item \path{tables/weather/tier1/prism_tier1_qa_findings.csv}
\item \path{tables/weather/tier1/prism_summary_by_county_review_flags.csv}
\item \path{tables/weather/tier1/prism_summary_pooled.csv}
\item \path{tables/weather/tier1/prism_interannual_variability_by_state.csv}
\item \path{tables/weather/tier1/prism_summary_by_calendar_month.csv}
\item \path{tables/weather/tier1/prism_tier1_coverage_integrity.csv}
\item \path{tables/weather/tier1/prism_summary_by_county.csv}
\item \path{tables/weather/tier1/prism_county_interannual_variability.csv}
\end{itemize}
\end{document}
